\documentclass{minmal}
\begin{document}
\noindent Skriv syre-base reaksjonen for følgende stoffer:\\
$H_2SO_4+2H_2O=SO_4^{2-}+2H_3O^+$\\
$HCl+H_2O=Cl^-+H_3O^+$\\
$HBr+H_2O=Br^-+H_3O^+$\\
$HClO_3+H_2O=ClO_3^-+H_3O^+$\\
$NaOH=Na^++OH^-$\\
$Ca(OH)_2=Ca^{2+}+2OH^-$\\\\
Bestem pH i 0,2M 0,1M og 0,3M løsnig for hver av stoffene over.\\\\
For å bestemme pH må vi huske definisjonen:\\
$pH=-log([H_3O^+])$\\
reaksjon 1:\\\\
$H_2SO_4+2H_2O=HSO_4^{2-}+2H_3O^+$  (syren protlyserer fullstendig)\\\\
Siden syren er sterk protolyserer den fullstendig i antar vi.\\\\

\noindent Det betyr at for $0,1M H_2SO_4$ så er $[H_3O^+]=2\cdot [H_2SO_4]=2\cdot0,1M=0,2M$\\\
$pH=-log([H_3O^+])=-log(0,2)=0,69$\\\\
0,2M $H_2SO_4$:
for $0,2M H_2SO_4$ så er $[H_3O^+]=2\cdot [H_2SO_4]=2\cdot0,2M=0,4M$\\\
$pH=-log([H_3O^+])=-log(0,4)=0,4$\\\\
0,3M $H_2SO_4$:
for $0,3M H_2SO_4$ så er $[H_3O^+]=2\cdot [H_2SO_4]=2\cdot0,3M=0,6M$\\\
$pH=-log([H_3O^+])=-log(0,6)=0,22$\\\\



\noindent Reaksjon 2:\\\\
$HCl+H_2O=Cl^-+H_3O^+$\\\\
$0,1M HCl: [H_3O^+]=[HCl]=0,1M$\\\\
$pH=-log([H_3O^+]=-log(0,1)=1$\\\\
$0,2M HCl: [H_3O^+]=[HCl]=0,2M$\\\\
$pH=-log([H_3O^+]=-log(0,2)=0,69$\\\\
$0,3M HCl: [H_3O^+]=[HCl]=0,3M$\\\\
$pH=-log([H_3O^+]=-log(0,3)=0,52$\\\\
\noindent Reaksjon 3:\\\\
$HBr+H_2O=Br^-+H_3O^+$\\\\
$0,1M HBr: [H_3O^+]=[HBr]=0,1M$\\\\
$pH=-log([H_3O^+]=-log(0,1)=1$\\\\
$0,2M HBr: [H_3O^+]=[HBr]=0,2M$\\\\
$pH=-log([H_3O^+]=-log(0,2)=0,69$\\\\
$0,3M HBr: [H_3O^+]=[HBr]=0,3M$\\\\
$pH=-log([H_3O^+]=-log(0,3)=0,52$\\\\

\noindent Reaksjon 4: Blir identisk med reaksjon 2 og 3.\\\\
\noindent Reaksjon 5:\\\\
$NaOH=Na^++OH^-$\\\\
0,1M NaOH: Her vil $[OH^-]=[NaOH]=0,1M$\\\
Dette er en base, vi må beregne pOH.\\\\
$pOH=-log([OH^-])=-log(0,1)=1$\\\\
$pH=14-POH=14-1=13$\\\\
0,2M NaOH: Her vil $[OH^-]=[NaOH]=0,2M$\\\
Dette er en base, vi må beregne pOH.\\\\
$pOH=-log([OH^-])=-log(0,2)=0,69$\\\\
$pH=14-POH=14-0,69=13,31$\\\\
0,3M NaOH: Her vil $[OH^-]=[NaOH]=0,3M$\\\
Dette er en base, vi må beregne pOH.\\\\
$pOH=-log([OH^-])=-log(0,3)=0,52$\\\\
$pH=14-POH=14-0,52=13,48$\\\\
\noindent Reaksjon 6:\\\\
$Ca(OH)_2=Ca^{2+}+2OH^-$\\\\
0,1M $Ca(OH)_2$: Her er $[OH^-]=2\cdot[Ca(OH)_2]=2\cdot 0,1M=0,2M$\\\\
Siden dette er en base, så må vi beregne pOH først:\\\\
$pOH=-log([OH^-]=-log(0,2)=0,69$\\\\
$pH=14-pOH=14-0,69=13,31$\\\\
0,2M $Ca(OH)_2$: Her er $[OH^-]=2\cdot[Ca(OH)_2]=2\cdot 0,2M=0,4M$\\\\
Siden dette er en base, så må vi beregne pOH først:\\\\
$pOH=-log([OH^-]=-log(0,4)=0,39$\\\\
$pH=14-pOH=14-0,49=13,61$\\\\
0,3M $Ca(OH)_2$: Her er $[OH^-]=2\cdot[Ca(OH)_2]=2\cdot 0,3M=0,6M$\\\\
Siden dette er en base, så må vi beregne pOH først:\\\\
$pOH=-log([OH^-]=-log(0,6)=0,22$\\\\
$pH=14-pOH=14-0,22=13,78$\\\\
\\\\Repetisjon fra i går:\\\\
Skriv syre-base reaksjonen for følgende stoffer:\\\\
$H_2SO_4+H_2O=HSO_4^-+H_3O^+$\\\\\\
$HCl+H_2O=Cl^-+H_3O^+$\\\\\\
$HBr+H_2O=Br^-+H_3O^+$\\\\\\
$HClO_3+H_2O=ClO_3^-+H_3O^+$\\\\\\
$NaOH=Na^++OH^-$\\\\\\
$Ca(OH)_2=Ca^{2+}+2OH^-$\\\\\\\\
$H_3PO_4+H_2O=H_3O^++H_2PO_4^-$\\\\
Skriv reaksjonen med vann for følgende stoffer:\\
$NH_3$ (er en base)\\
Alle tre trinn for fosforsyre ($H_3PO_4$)\\
$HNO_3$ syre\\
$CaO$ base\\
Øvingsoppagver med pH:
\noindent
Bestem pH i 0,1M HCl, 0,2M NaOH, 0,2M $Ca(OH)_2$ 0,05M $H_2SO_4$ som protolyserer fullstendig. \\
Du skal også skrive syre-base reaksjonen for hver av dem (reaksjon med vann).\\\\
Vi har 20mL 0,1M HCl som vi titrer med NaOH. Det går med 15mL $NaOH$.  Hva er konsentrasjonen av $NaOH$? \\\\
\noindent
Hva blir konsentrasjonen av basen (her er det viktig med riktig reaksjon)?\\\\
1) 25mL 0,2M HCl titreres mot 15mL NaOH?\\\\
2) 15mL 0,1M HCl titreres mot 15mL NaOH  (ikke lov med kalkulator på denne)?\\\\
3) 25mL 0,3M HCl titreres mot 20mL NaOH?\\\\
4) 25mL 0,2M HCl titreres mot 15mL $Ca(OH)_2$? 
\end{document}